\makeatletter
\def\input@path{{../styles/}{../../styles/}{../../../styles/}{../}{../../}{../../../}}
\makeatother

\documentclass{ee102_pset}
% macros.tex - Course meta information
\renewcommand{\course}{EE 102} % with a space
\renewcommand{\coursetitle}{Signal Processing and Linear Systems}
\renewcommand{\instructor}{Ayush Pandey}
\renewcommand{\student}{Name: }

\renewcommand{\semester}{Fall 2025}
\date{\semester} % this sets the LaTeX date field safely

% Problem set number
\renewcommand{\psetnum}{1}

% Release / due — use renewcommand because package provides empties
\renewcommand{\releasedate}{August 20, 2025}
\renewcommand{\duedate}{August 27, 2025}

\input{preamble.tex}

\author{\rule{3cm}{0.4pt}}
\submitdate{\rule{3cm}{0.4pt}}
\problemset{\shortstack{Problem Set \#9: Fourier Transform Properties\\and RC Applications}}
\renewcommand{\psetnum}{9}
\renewcommand{\releasedate}{October 26, 2026}
\renewcommand{\duedate}{Nov 4, 2026 at 11:59 PM}
\shorttitle{Problem Set \#9}
\hypersetup{pdftitle={EE 102 Problem Set 9: Fourier Transform Properties and RC Applications}}

\setlength{\parskip}{8pt plus 1pt minus 1pt}
\setlength{\abovedisplayskip}{6pt plus 2pt minus 2pt}
\setlength{\belowdisplayskip}{6pt plus 2pt minus 2pt}

\begin{document}
Released October 26, 2026. Due \textbf{Wednesday, November 4, at 11:59 PM}
(Pacific time). This shorter exam-week assignment has two written problems
and is worth \textbf{100 points}. No programming is required.

Use the continuous-time Fourier transform convention
\[
X(\omega)=\int_{-\infty}^{\infty}x(t)e^{-j\omega t}\,dt,
\qquad
x(t)=\frac{1}{2\pi}\int_{-\infty}^{\infty}X(\omega)e^{j\omega t}\,d\omega.
\]
Frequencies $\omega$ are in rad/s. Show your reasoning; exact expressions
are sufficient.

\problem{1} The RC circuit in the frequency domain. \hfill [50 points]
% Adapted from Fall 2025 HW8 Problem 3, with the earlier LTI eigenfunction application.

A series RC circuit has $R=1\,\mathrm{k}\Omega$ and
$C=1000\,\mu\mathrm{F}$. The input $x(t)$ is the voltage across the series
combination; the output $y(t)$ is the capacitor voltage. With zero initial
stored energy, its impulse response is
\[
h_C(t)=\frac{1}{\tau_C}e^{-t/\tau_C}u(t),\qquad \tau_C=RC.
\]

\problempart[15 points] Calculate $\tau_C$. Derive $H_C(\omega)$ directly
from the Fourier transform integral, and find $H_C(0)$.

\problempart[15 points] Find the magnitude and phase of $H_C(\omega)$ at
$\omega=1$ and $\omega=3$ rad/s. What happens to the magnitude as
$|\omega|\to\infty$?

\problempart[20 points] The input is
\[
x(t)=2\cos(t)+2\cos(3t)\quad\text{volts},
\]
present for all time. Find the bounded steady-state output voltage. Which
input component is reduced more, and how does that relate to the name
``low-pass filter''?

Hint: Use the LTI complex-exponential eigenfunction property from Weeks 6--7.
A component $A\cos(\omega t)$ produces
$A|H_C(\omega)|\cos(\omega t+\angle H_C(\omega))$ in steady state.

\problem{2} Symmetry and time shifting of a pulse. \hfill [50 points]

Define the aperiodic pulse
\[
p(t)=
\begin{cases}
2,&-1<t<1,\\
0,&\text{otherwise}.
\end{cases}
\]

\problempart[15 points] Sketch $p(t)$. Before calculating its transform,
use its symmetry to predict whether $P(\omega)$ is real and whether it is
even. Then derive $P(\omega)$ and give its value at $\omega=0$.

Hint: For a real, even signal,
$P(\omega)=2\int_0^{\infty}p(t)\cos(\omega t)\,dt$.
At zero frequency, use the signed area of the pulse.

\problempart[20 points] Let $q(t)=p(t-2)$. Sketch $q(t)$ and obtain
$Q(\omega)$ using the time-shifting property. Explain why
$|Q(\omega)|=|P(\omega)|$ even though $q(t)$ is not even about $t=0$.
Find the magnitude and phase of $Q(\pi/4)$.

Hint: $x(t-t_0)\longleftrightarrow e^{-j\omega t_0}X(\omega)$,
$\sin(\pi/4)=\sqrt{2}/2$, and $e^{-j\pi/2}=-j$.

\problempart[15 points] Now let $r(t)=p(t-2)+p(t+2)$. Without performing
another transform integral, find $R(\omega)$ and write it as a real,
even expression. Evaluate $R(0)$ and $R(\pi/4)$. Explain how the two
shifted pulses can cancel at one frequency even though both pulses
are positive in time.

Hint: Apply linearity and time shifting, then use
$e^{j\theta}+e^{-j\theta}=2\cos\theta$.

\end{document}
